\documentclass[11pt]{article}
\usepackage[margin=1in]{geometry}
\usepackage{booktabs}
\usepackage{threeparttable}
\usepackage{amsmath}
\usepackage{array}
\usepackage{microtype}
\usepackage{enumitem}
\usepackage{hyperref}
\usepackage{setspace}
\usepackage{graphicx}
\setlength{\parindent}{0pt}
\setlength{\parskip}{0.65em}
\hypersetup{colorlinks=true,linkcolor=black,urlcolor=black,citecolor=black}

\newcommand{\memoTable}[2]{%
  \IfFileExists{#1}{\input{#1}}{%
    \begin{table}[!htbp]\centering\small
    \caption{#2}
    \fbox{\parbox{0.90\textwidth}{\centering
    This exact five-column regression table is generated by \\texttt{R/13\_generate\_coauthor\_memo\_tables.R}. No coefficient is inserted here until it has been re-estimated from the final analysis data and Model~4 has reproduced the frozen preferred-model coefficient.}}
    \end{table}%
  }%
}

\begin{document}

\begin{center}
{\Large\bfseries Memo: Recent Empirical Results and Identification Work}\\[0.25em]
{\normalsize August 10, 2026}
\end{center}

Over the last several days, I have substantially expanded the empirical analysis of the relationship between foreign direct investment (FDI), demographic context, centrism, and BJP support. The emerging picture is increasingly coherent across constituency- and individual-level analyses. At the same time, extensive work on quasi-experimental identification suggests that our strongest claims should presently remain longitudinal and associational rather than fully causal.

\section*{Headline substantive results}

The most consistent evidence concerns \textbf{manufacturing FDI interacting with pre-existing demographic context}. Constituencies experiencing greater manufacturing FDI between the 2009 and 2014 elections tend to shift further toward the BJP where Muslim presence or established migration was already greater. These relationships are much weaker for services FDI and for measures of recent demographic change. The pattern is therefore more consistent with new economic exposure interacting with an established social environment than with globalization and demographic change simply occurring simultaneously.

\subsection*{How to read the percentage-point contrasts}

\textbf{The percentage-point quantities in Table~\ref{tab:bridge} are not regression coefficients.} They are post-estimation substantive contrasts that put the aggregate and respondent models on a common, interpretable geographic ruler. The underlying regressions are estimated on their own samples and in their native units; no model is re-estimated on a common sample. Only the numeric movements used to translate the fitted interaction coefficient into an interpretable quantity are standardized.

For the \textbf{Muslim-share two-way contrast}, manufacturing FDI is moved from zero to the median positive local exposure, $0\rightarrow0.5346$ in $\log(1+\text{projects per 100,000})$, while 2001 Muslim population share is moved from the 25th to the 75th percentile, $3.72\%\rightarrow15.07\%$. The reported number answers: \emph{How much larger is the predicted effect of that FDI increase in a high-Muslim-share constituency than in a low-Muslim-share constituency?} Thus, an estimate of $+1.50$ means that the specified increase in FDI is associated with an approximately 1.50-percentage-point larger BJP shift at the 75th percentile of Muslim share than at the 25th percentile, not that a one-unit increase in either raw regressor raises BJP support by 1.50 points.

For the \textbf{established-migration two-way contrast}, the analogous FDI movement is $0\rightarrow0.5273$ and the established migrant-stock share moves from $0.773\%\rightarrow3.827\%$. The migration moderator is \texttt{mig\_total\_upto\_2001\_share\_ac\_pop}: interstate plus international migrants who report 10--19 or 20+ years of residence in the 2011 Census, scaled by AC population. It therefore proxies the migrant stock already resident by roughly 2001; it is not a literal 2001 Census migration-share measure.

The \textbf{three-way contrasts answer a second-difference question}. At the aggregate level, they ask how much the FDI-by-demographic contrast itself changes as the survey-weighted 2009 constituency Center share moves from its 25th to its 75th percentile ($0\rightarrow9.07\%$ in the common reference distribution). At the respondent level, the Center contrast instead moves an individual from non-Center to Center ($0\rightarrow1$) among ideology-complete respondents. Accordingly, aggregate and respondent triple-interaction magnitudes are substantively related but \textbf{must not be mechanically compared}: one moderates by a constituency composition, the other by an individual characteristic. The regression tables in the appendix report the \emph{raw fitted coefficients}, standard errors, and p-values; Table~\ref{tab:bridge} reports only these transformed substantive contrasts.

\begin{table}[!htbp]
\centering
\small
\caption{Preferred cross-level substantive interaction contrasts}
\label{tab:bridge}
\begin{threeparttable}
\begin{tabular}{llccc}
\toprule
Demographic context & Interaction & First difference & Lagged outcome & Respondent, 2014 \\
\midrule
Muslim share, 2001 & FDI $\times$ demographic & 1.50 [0.36, 2.63] & 1.87 [0.71, 3.02] & 2.84 [0.04, 5.63] \\
Muslim share, 2001 & FDI $\times$ demographic $\times$ Center & 4.21 [1.08, 7.34] & 3.77 [0.63, 6.91] & 2.55 [$-4.72$, 9.81] \\
Established migration & FDI $\times$ demographic & 0.41 [$-0.39$, 1.21] & 0.42 [$-0.54$, 1.37] & 3.94 [2.14, 5.75] \\
Established migration & FDI $\times$ demographic $\times$ Center & 0.39 [$-0.79$, 1.57] & 0.57 [$-0.84$, 1.97] & $-0.99$ [$-6.12$, 4.14] \\
\bottomrule
\end{tabular}
\begin{tablenotes}[flushleft]
\scriptsize
\item Notes: Entries are percentage-point substantive contrasts with 95\% confidence intervals, not raw regression coefficients. FDI is manufacturing, local (own AC plus touching ACs), all announced/opened projects, log-transformed per 100,000 residents. The aggregate models use C1 controls. Respondent estimates use the weighted 2014 baseline-adjusted LPM with V2+C1, the candidate-present sample, state fixed effects, and PC-by-district multiway clustered standard errors. See the text immediately above for the exact contrast definition.
\end{tablenotes}
\end{threeparttable}
\end{table}

The clearest cross-level replication is therefore the \textbf{Muslim-context two-way relationship}: all three preferred estimates are positive and their 95\% intervals exclude zero. The evidence on migration differs across levels: the aggregate interaction is small and imprecise, whereas the respondent-level association is sizeable and precisely estimated. The evidence for centrism is also more nuanced than a simple ``centrists respond more'' story. Constituencies with larger pre-existing centrist electorates show a substantially stronger FDI-by-Muslim relationship---roughly a four-percentage-point amplification in both dynamic constituency designs---but the corresponding individual Center interaction is positive and imprecise. The safer interpretation is that \textbf{centrist-heavy constituencies appear to be especially responsive political environments, but the aggregate contextual result cannot yet be attributed uniquely to stronger behavioral responsiveness among individual centrist voters}.

The NES analysis is important because it indicates that the aggregate pattern corresponds to voter choice rather than only electoral arithmetic. The preferred respondent models use the 2014 election, condition on 2009 BJP support and matched pre-period FDI, apply survey weights, and restrict the primary sample to voters for whom a BJP candidate was available. The Muslim-context result is robust across weighted LPM, logit, unweighted, and broader voter-sample specifications. The migration-context relationship is especially strong at the respondent level.

We separately tested whether the findings could be generated by \textbf{BJP candidate supply} rather than voter demand. Candidate availability itself shows little evidence of responding to FDI-by-demographic context. Restricting the individual analysis to parliamentary constituencies in which the BJP contested in both 2009 and 2014 also leaves the principal two-way effects largely intact: the migration contrast is 5.38 percentage points in the original candidate-present analysis and 5.40 among always-contested constituencies; the Muslim contrast is 3.35 versus 2.75 points, respectively, with the latter less precise. Strategic candidate entry therefore does not appear to explain the core voter-level relationships.

\section*{Why the headline tables do not report pooled 2009/2014 models}

We estimated pooled specifications and retain them as robustness checks, but they are not the headline estimands and are intentionally not reproduced in the tables below. At the constituency level, pooling two election years combines persistent between-constituency differences with within-constituency temporal change. That is a particularly awkward estimand here because the core demographic moderators are pre-existing stocks and many contextual measures are fixed or nearly fixed across the two election rows. By contrast, the first-difference model directly asks which constituencies shifted more toward the BJP between 2009 and 2014 as a function of 2009--14 FDI exposure, while the lagged-outcome model asks which constituencies had greater BJP support in 2014 conditional on their 2009 BJP support and pre-period FDI. These dynamic designs therefore align much more closely with the temporal argument of the paper.

The same logic is stronger at the respondent level. The NES data are repeated cross-sections rather than a panel of the same voters, so a pooled 2009/2014 individual model does not identify within-person switching and combines two different electoral environments. Our preferred respondent design instead uses 2014 vote choice, measures FDI over the intervening 2009--14 window, and conditions on the constituency's 2009 FDI exposure and BJP vote share. Pooled results are therefore useful specification checks, but presenting them beside the dynamic models as if they estimated the same quantity would obscure rather than clarify the main empirical comparison.

\section*{Identification and shift-share work}

Before examining any political 2SLS coefficient, we set four pre-outcome credibility requirements for a candidate instrument:

\begin{enumerate}[label=(\roman*),leftmargin=2.3em]
\item \textbf{First-stage relevance.} The instrument should meaningfully predict 2009--14 manufacturing FDI after predetermined controls. We report the clustered single-instrument first-stage $F$ and partial $R^2$; $F\approx10$ is a useful screening convention, not a mechanical proof of strength.
\item \textbf{A clean pre-period placebo.} The same putatively future external shock should not strongly predict 2004--09 FDI after the corresponding baseline controls. A strong placebo implies that the instrument is tracking persistent FDI geography or industrial selection rather than generating genuinely new post-2009 exposure.
\item \textbf{Stability of the first stage.} Relevance should not disappear when one major industry or source country is removed. Leave-one-out diagnostics are especially important for shift-share designs because an apparently strong aggregate first stage can in fact be generated by one or two influential shocks.
\item \textbf{Diffuse, plausibly external shock variation and a defensible exclusion restriction.} Identifying weight should be spread across a reasonable number of shocks, and the shocks should affect BJP support primarily through local FDI rather than through other channels such as export demand, domestic employment, prices, or sector-specific political change.
\end{enumerate}

Against those standards, we tested four related strategies:

\begin{enumerate}[label=\arabic*.,leftmargin=2.1em]
\item \textbf{Global industry Bartik / within-manufacturing Bartik.} We combined 2005 Economic Census industrial composition with subsequent worldwide greenfield-FDI industry shocks from WIR/fDi Markets. The most defensible version conditions on total manufacturing share and uses only within-manufacturing composition, which prevents the instrument from simply distinguishing manufacturing-heavy from non-manufacturing ACs. For own-AC manufacturing FDI, the all-complete EC05 sample gives $F_{\text{post}}=3.60$ and $F_{\text{placebo}}=1.12$; restricting to high-quality EC05 observations raises the post-period first stage to $F_{\text{post}}=10.25$ with partial $R^2=0.0044$, while the placebo remains low at $F_{\text{placebo}}=1.11$. This is the cleanest relevance/placebo combination we found. The weakness is concentration: the effective number of industry shocks is only about 3.4. In the all-complete leave-one-industry-out audit, the post-period $F$ falls from 3.60 to 1.07 when textiles/clothing/leather is removed, 2.07 when food/beverages/tobacco is removed, and 2.77 when non-metallic minerals is removed. The high-quality headline result therefore does not establish a broadly supported first stage; much of the identifying variation rests on a small set of sectors. That makes both weak-instrument sensitivity and the exclusion restriction difficult to defend, because large global shocks to those sectors can plausibly influence local politics through trade, employment, and prices even absent local FDI.

\item \textbf{Matched-local Bartik.} Because the observational treatment is local FDI (own AC plus touching ACs), we built a geographically matched instrument using the same local economy on the exposure side. In the strict complete-neighborhood EC05 sample the within-manufacturing local instrument gives $F_{\text{post}}=3.74$ but $F_{\text{placebo}}=10.30$. Restricting to neighborhoods whose EC05 members are all high quality raises the post-period first stage to $F_{\text{post}}=12.66$, but the placebo remains substantial at $F_{\text{placebo}}=7.79$. Thus the local version can be relevant in a selected high-quality sample, but it also strongly predicts where FDI was already located before the treatment window---exactly the pattern we do not want from a clean external instrument.

\item \textbf{Official 2012--13 Indian FDI liberalizations.} We audited the sector-specific reforms as potential policy instruments and cross-walked them to both EC05/SHRIC employment exposure and the fDi Markets project classification. Three reform areas have clean national EC05 exposure mappings---scheduled/non-scheduled air transport, courier services, and telecom services---but \textbf{none cleanly targets the paper's preferred manufacturing-FDI treatment}. Only one policy sector also maps cleanly to the available fDi Markets project-sector taxonomy. The reform package therefore does not provide a defensible national instrument for manufacturing FDI during the paper's treatment window without redefining the treatment around a narrow set of service sectors.

\item \textbf{Source-country-push Bartik.} We matched source country to approximately 97.7\% of the geocoded Indian FDI projects and combined each industry's pre-2009 investor-country composition with subsequent outward greenfield-FDI shocks from those source countries. This materially improves shock diversification relative to the industry-only Bartik, but it fails the placebo gate. Using approximate rest-of-world source shocks, the all-EC05 sample gives $F_{\text{post}}=3.78$ and $F_{\text{placebo}}=9.72$; the high-quality sample gives $F_{\text{post}}=7.19$ and $F_{\text{placebo}}=9.57$. The world-inclusive variant performs similarly or worse: $2.84$ versus $9.73$ in the full sample and $5.47$ versus $10.38$ in the high-quality sample. The source shocks therefore predict historical FDI geography at least as strongly as new 2009--14 FDI, indicating that the instrument is still carrying persistent investor-location patterns rather than a clean new capital-supply shock.
\end{enumerate}

The important methodological point is that these tests were adjudicated \textbf{before looking at any BJP IV second stage}. We therefore have not selected an instrument because it produces an attractive political coefficient. The current quasi-experimental evidence does not meet our pre-specified standard for clean instrumentation, so no BJP 2SLS result is treated as a headline causal estimate.

\section*{Bottom line}

The empirical case has become stronger but also better bounded. \textbf{Manufacturing FDI is robustly associated with increased BJP support in constituencies characterized by pre-existing demographic diversity; the Muslim-context relationship replicates particularly clearly across electoral returns and individual vote choice; and the main voter-level findings are not readily explained by BJP candidate supply.} Centrist \emph{contexts} appear especially important for the Muslim relationship, although the evidence for uniquely stronger responsiveness among individual centrist voters is much less certain.

\clearpage
\appendix
\section*{Publication-format regression tables}

The following tables report the raw fitted regression coefficients corresponding to the preferred Muslim and established-migration families. The five columns are deliberately cumulative. For aggregate models, Model~1 includes the interaction and design-defining baseline terms only; Model~2 adds C0; Model~3 adds C1; Model~4 is the preferred C1 specification with state fixed effects and PC-clustered standard errors; Model~5 retains Model~4 and adds the full C3 block (which nests C2). For respondent LPMs, Model~1 includes the interaction plus 2009 FDI and 2009 BJP support; Model~2 adds C0; Model~3 adds V2 and C1; Model~4 adds state fixed effects and PC-by-district multiway clustering and is the preferred model; Model~5 adds C3 while retaining V2, state fixed effects, clustering, survey weights, and the candidate-present sample. Cells report the coefficient, standard error in parentheses, and exact p-value.

\memoTable{data/derived/switchers_rewrite/model_exploration/coauthor_empirical_memo/tables/aggregate_muslim_two_way_first_difference.tex}{Muslim share, two-way first-difference model}
\memoTable{data/derived/switchers_rewrite/model_exploration/coauthor_empirical_memo/tables/aggregate_muslim_two_way_lagged_outcome.tex}{Muslim share, two-way lagged-outcome model}
\memoTable{data/derived/switchers_rewrite/model_exploration/coauthor_empirical_memo/tables/aggregate_muslim_triple_first_difference.tex}{Muslim share, three-way first-difference model}
\memoTable{data/derived/switchers_rewrite/model_exploration/coauthor_empirical_memo/tables/aggregate_muslim_triple_lagged_outcome.tex}{Muslim share, three-way lagged-outcome model}
\memoTable{data/derived/switchers_rewrite/model_exploration/coauthor_empirical_memo/tables/aggregate_migration_two_way_first_difference.tex}{Established migration, two-way first-difference model}
\memoTable{data/derived/switchers_rewrite/model_exploration/coauthor_empirical_memo/tables/aggregate_migration_two_way_lagged_outcome.tex}{Established migration, two-way lagged-outcome model}
\memoTable{data/derived/switchers_rewrite/model_exploration/coauthor_empirical_memo/tables/aggregate_migration_triple_first_difference.tex}{Established migration, three-way first-difference model}
\memoTable{data/derived/switchers_rewrite/model_exploration/coauthor_empirical_memo/tables/aggregate_migration_triple_lagged_outcome.tex}{Established migration, three-way lagged-outcome model}

\memoTable{data/derived/switchers_rewrite/model_exploration/coauthor_empirical_memo/tables/respondent_muslim_two_way.tex}{Respondent weighted LPM: Muslim share, two-way interaction}
\memoTable{data/derived/switchers_rewrite/model_exploration/coauthor_empirical_memo/tables/respondent_muslim_triple.tex}{Respondent weighted LPM: Muslim share, three-way interaction}
\memoTable{data/derived/switchers_rewrite/model_exploration/coauthor_empirical_memo/tables/respondent_migration_two_way.tex}{Respondent weighted LPM: established migration, two-way interaction}
\memoTable{data/derived/switchers_rewrite/model_exploration/coauthor_empirical_memo/tables/respondent_migration_triple.tex}{Respondent weighted LPM: established migration, three-way interaction}

\end{document}
